\documentclass[10pt,a4paper]{amsart}
\usepackage[margin=19mm]{geometry}
\usepackage{amsmath,amssymb,mathtools}
\usepackage[scr=rsfs,cal=euler]{mathalfa}
\usepackage{xfrac}
\usepackage[unicode,hidelinks,bookmarks=false]{hyperref}
\newcommand{\ie}{i.e.\ }
\newcommand{\cf}{cf.\ }
\newcommand{\resp}{respectively\ }
\newcommand{\Subsection}{\subsection}
\providecommand{\nobreakdash}{\nobreak\hbox{-}}
\setlength{\emergencystretch}{2em}
\multlinegap0pt
\usepackage{amssymb}
\usepackage[shortlabels]{enumitem}
\newtheorem{lemma}{Lemma}
\newtheorem{theorem}{Theorem}
\title{Hermitian Distance Degree of Unitary-Invariant Matrix Varieties}
\author{Nikhil Ken}
\date{Source: arXiv:2602.10737v2 (CC BY 4.0)}
\begin{document}
\maketitle

\noindent\emph{Source and licence.} Hermitian Distance Degree of Unitary-Invariant Matrix Varieties. Version arXiv:2602.10737v2 (CC BY 4.0), distributed by arXiv under the Creative Commons Attribution 4.0 International licence, http://creativecommons.org/licenses/by/4.0/, as recorded in the arXiv licence metadata for this exact version and shown on the abstract page https://arxiv.org/abs/2602.10737v2. Full licence notice: \texttt{licenses/arxiv-2602.10737-CC-BY-4.0.txt}.\par\smallskip

\noindent\textit{Editorial scope.} Counted unit: the article's theorem gg (source0.tex lines 233--257). The article's complete original proof of it is delivered with it (source0.tex lines 285--365, one \texttt{proof} environment of 81 lines, which the article places away from the statement). No mathematics is written, repaired or re-derived: the statement and the proof are the article's own lines, in the article's own notation, and no retained line is substituted or re-flowed.  Retained uncounted as the source-local context the proof needs: lines 190--195.  Nothing was omitted from any retained span, so the package carries no omission record.  No retained line contains a citation, so the wrapper carries no bibliography and no bibliography entry.  No retained line contains a figure environment, an image-inclusion command or drawing code, so no image is delivered and the package declares no asset dependency.  Editorial formatting only, in addition: an AMS \texttt{amsart} preamble that restates the article's own declarations and macros used by the retained lines, namely the environments \texttt{lemma}, \texttt{theorem} on separate counters, together with the standard packages those lines need.  The retained environments print in this document as Lemma 1, Theorem 1 in order of appearance, whereas the article prints the same items with the numbering of its own full document.  The complete native source of the article as distributed in the arXiv e-print for this exact version is bundled unchanged as \texttt{sources/194424/source0.tex}.
\begin{lemma}\label{a}
   If $A_1 = u_1 v_1^*$ and $A_2 = u_2 v_2^*$ are rank-one matrices, then the Hermitian inner product given by $\langle A_1, A_2 \rangle = \mathrm{Tr}(A_1 A_2^*)$ is also given by

$$
\langle A_1, A_2 \rangle = \langle u_1, u_2 \rangle \langle v_1, v_2 \rangle.
$$\end{lemma}

\begin{theorem}[\textbf{Eckart-Young for HD}]\label{gg}
Let $A \in \mathbb{C}^{m \times n}$ be a matrix of rank $r$, and let its singular value decomposition be
$$
A = U \Sigma V^*, \qquad 
\Sigma = \mathrm{diag}(s_1, \dots, s_r, 0, \dots, 0),
$$
with singular values $s_1 \geq \cdots \geq s_r > 0$.  

For $1 \leq i \leq r$, define 
$$
\Sigma_i = \mathrm{diag}(0, \dots, 0, s_i, 0, \dots, 0),
$$
so that
$$
A = U \Sigma_1 V^* + U \Sigma_2 V^* + \cdots + U \Sigma_r V^*.
$$

Then the critical points of the Hermitian distance  function from $A$ to the variety of matrices of rank at most $k$ are precisely the matrices of the form
$$
U \bigl(\Sigma_{i_1} + \cdots + \Sigma_{i_k}\bigr) V^*, 
\qquad 1 \leq i_1 < \cdots < i_k \leq r.
$$

In particular, if the nonzero singular values of $A$ are distinct, then for each $k \leq r$ there are exactly $\binom{r}{k}$ such critical points.
\end{theorem}

\begin{proof}[Proof of Theorem~\ref{gg}]

 The matrix \( U(\Sigma_{i_1} + \cdots + \Sigma_{i_k}) V^* \) is a critical point of the distance function from \( A \) to the variety \( M_k \) if and only if the vector 

\[
A - U(\Sigma_{i_1} + \cdots + \Sigma_{i_k}) V^*
\]

is orthogonal to the tangent space.  \[
T_{U(\Sigma_{i_1} + \cdots + \Sigma_{i_k}) V^*} M_k = (\mathbb{C}^m \otimes{ v^*_{i_1}}+ u_{i_1} \otimes \mathbb{C}^n) + \cdots + (\mathbb{C}^m \otimes {v^*_{i_k}} + u_{i_k} \otimes \mathbb{C}^n).
\]

From the SVD of \( A \), we have 

\[
A - (U(\Sigma_{i_1} + \cdots + \Sigma_{i_k}) V^*) = U(\Sigma_{j_1} + \cdots + \Sigma_{j_l}) V^* = s_{j_1} u_{j_1} \otimes v^*_{j_1} + \cdots + s_{j_l} u_{j_l} \otimes v^*_{j_l}
\]

where \( \{j_1, \ldots, j_l\} \) is the set of indices given by the difference \( \{1, \ldots, r\} \setminus \{i_1, \ldots, i_k\} \). Let $e_1,\dots,e_m$ be the canonical basis of $\mathbb{C}^m$, then by Lemma \ref{a} We have \[
\langle s_{j_h} u_{j_h} \otimes v^*_{j_h}, e_l \otimes v^*_{i_s} \rangle = s_{j_h} \langle u_{j_h}, e_l \rangle \langle v_{j_h}, v_{i_s} \rangle = 0
\]
since \( v_{j_h} \) and \( v_{i_s} \) are distinct columns of the unitary matrix \( V \). Thus, the matrices \( U\Sigma_{j_h} V^* \) are orthogonal to the spaces \( \mathbb{C}^m \otimes V^*_{i_s} \).

In a similar way, since \( U \) is a unitary matrix, the matrices \( U\Sigma_{j_h} V^* \) are orthogonal to the spaces \( u_{i_s} \otimes \mathbb{C}^n \). Therefore, \( A - (U(\Sigma_{i_1} + \cdots + \Sigma_{i_k}) V^*) \) is orthogonal to the tangent space, and \( U(\Sigma_{i_1} + \cdots + \Sigma_{i_k}) V^* \) is a critical point. Now we will show all critical points are of this form, Let $B\in M_k$ be such that $A-B$ is orthogonal to the tangent space $T_B(M_k)$. We consider the SVD of B and A-B,let \( B = U'(\Sigma_1' + \cdots + \Sigma_k')V'^* \), and $
A - B = U''(\Sigma_1'' + \cdots + \Sigma_l'')V''^*
$
with \( \Sigma'_k \neq 0 \) and \( \Sigma''_l \neq 0 \). Since A-B is orthogonal to $T_B(M_k)$ we get that $\langle Col(A-B),u_s'\rangle=0=\langle Row(A-B),v_s'\rangle$ for $s= 1 ,\dots, k$. In particular, \(\text{Col}(A - B)\) is a vector subspace of \(\text{Span}\{u'_1, \ldots, u'_k\}^\perp\) and has dimension at most \(m - k\), while \(\text{Row}(A - B)\) is a vector subspace of \(\text{Span}\{\overline{v}'_1, \ldots, \overline{v}'_k\}^\perp\) and has dimension at most \(n - k\), so that 

$
l \leq \min\{m, n\} - k.
$
From the equality 

\[
A - B = (u''_1, \dots, u''_l, 0, \dots, 0)(\Sigma''_1 + \cdots + \Sigma''_l)V''^*
\]

we get 

\[
\text{Col}(A - B) \subset \text{Span}\{u''_1, \dots, u''_l\}
\]

and equality holds because the rank of A-B is l. In a similar way, 

\[
\text{Row}(A - B) = \text{Span}\{\overline{v}''_1, \dots, \overline{v}''_l\}.
\]

This implies that the orthonormal columns \( u''_1, \dots, u''_l, u'_1, \dots, u'_k \) can be completed with orthonormal \( m - l - k \) columns of \( \mathbb{C}^m \) to obtain a unitary \( m \times m \) matrix \( U \), while the orthonormal columns \( v''_1, \dots, v''_l, v'_1, \dots, v'_k \) can be completed with orthonormal \( n - l - k \) columns of \( \mathbb{C}^n \) to obtain a unitary \( n \times n \) matrix \( V \). We get 

\[
A - B = U
\begin{pmatrix}
\Sigma'' & 0 & 0 \\
0 & 0 & 0 \\
0 & 0 & 0
\end{pmatrix}
V^*,
\]

\[
B = U
\begin{pmatrix}
0 & 0 & 0 \\
0 & \Sigma' & 0 \\
0 & 0 & 0
\end{pmatrix}
V^*,
\]

where 

\[
\Sigma'' = \mathrm{diag}(s''_1, \dots, s''_l), \quad \Sigma' = \mathrm{diag}(s'_1, \dots, s'_k).
\]
 Thus $B$ is obtained by keeping exactly $k$ of the rank one summands $U\Sigma_iV^*$ in the SVD expansion of $A$, hence it has the required form.



\end{proof}

\end{document}
